\chapter{The Single-Block Fundamental Theorem}\label{ch:single}

An injective MPS tensor that generates the same matrix-product vectors as a
second tensor of equal bond dimension is conjugate to
it~\cite[Corollary~IV.5]{Cirac2021Matrix}: there is an invertible~$X$ with
$B^i = X A^i X^{-1}$ for every physical index~$i$. The argument is purely
algebraic. Equality of the matrix-product vectors gives $\tr(A^w)=\tr(B^w)$ for
every word~$w$; the length-two and length-three traces produce a linear map~$T$
with $T(A^i)=B^i$ that is multiplicative and nonzero; a nonzero multiplicative
endomorphism of the simple algebra~$\MN{D}$ is a unital automorphism; and by
Skolem--Noether every automorphism of~$\MN{D}$ is conjugation by an
invertible~$X$. Evaluating at~$A^i$ closes the proof. The same conjugation, one
normal block at a time, is the central step of the multi-block theorem in
Chapter~\ref{ch:ft_proof}.

\section{The multiplicative linear extension}

\begin{theorem}[Nondegeneracy of the trace pairing]\label{thm:trace_nondeg}
    \lean{Matrix.trace_mul_right_eq_zero_iff}
    \leanok
    For $M \in \MN{D}$, $\tr(M N) = 0$ for all $N \in \MN{D}$ if and only if
    $M = 0$.
\end{theorem}

\begin{proof}
    \leanok
    Testing against the matrix unit $E_{ji}$ gives $\tr(E_{ji} M) = M_{ij}$, so
    $\tr(M N) = 0$ for all~$N$ forces every entry of~$M$ to vanish.
\end{proof}

\begin{lemma}[Same matrix-product vectors give trace agreement]\label{lem:same_mpv_trace}
    \lean{MPSTensor.SameMPV.trace_evalWord}
    \leanok
    \uses{def:same_mpv, def:eval_word}
    If $A$ and $B$ generate the same matrix-product vectors, then
    $\tr(A^w) = \tr(B^w)$ for every word~$w$.
\end{lemma}

\begin{proof}
    \leanok
    \uses{def:mpv}
    The configuration $\sigma = w$ gives $V^{(|w|)}(A)_w = \tr(A^w)$. Equality of
    the matrix-product vectors gives $V^{(N)}(A) = V^{(N)}(B)$ for all~$N$, so
    $\tr(A^w) = \tr(B^w)$.
\end{proof}

The trace conditions constrain every product of the matrices~$A^i$. To bring the
algebra structure of~$\MN{D}$ to bear, record how a matrix pairs against the
family $\{A^i\}$ under the trace.

\begin{definition}[Trace pairing map]\label{def:trace_pairing_map}
    \lean{MPSTensor.traceMulRightPi}
    \leanok
    \uses{def:mps_tensor}
    For a tensor~$A$, the trace pairing map $\Phi_A : \MN{D} \to \C^d$ is
    \[
        \Phi_A(M)_i = \tr(M A^i).
    \]
\end{definition}

\begin{theorem}[The trace pairing map is injective for injective tensors]%
    \label{thm:trace_pairing_injective}
    \lean{MPSTensor.traceMulRightPi_ker_eq_bot}
    \leanok
    \uses{def:trace_pairing_map, def:injective}
    If $A$ is injective, then $\ker \Phi_A = \{0\}$.
\end{theorem}

\begin{proof}
    \leanok
    \uses{thm:trace_nondeg}
    If $\Phi_A(M) = 0$, then $\tr(M A^i) = 0$ for every~$i$. The $\{A^i\}$ span
    $\MN{D}$, so $\tr(M N) = 0$ for all~$N$, and $M = 0$ by
    Theorem~\ref{thm:trace_nondeg}.
\end{proof}

\begin{lemma}[Injectivity transfers across a range inclusion]%
    \label{lem:injectivity_transfer}
    \lean{MPSTensor.ker_bot_of_range_le}
    \leanok
    Let $\Phi, \Psi : V \to W$ be $\C$-linear maps with $V$ finite-dimensional.
    If $\ker \Phi = \{0\}$ and $\operatorname{range} \Phi \subseteq
        \operatorname{range} \Psi$, then $\ker \Psi = \{0\}$.
\end{lemma}

\begin{proof}
    \leanok
    By rank--nullity, $\ker \Phi = \{0\}$ gives $\dim \operatorname{range} \Phi =
        \dim V$, and the range inclusion gives
    \[
        \dim V = \dim \operatorname{range} \Phi
        \le \dim \operatorname{range} \Psi \le \dim V,
    \]
    so $\dim \operatorname{range} \Psi = \dim V$ and $\ker \Psi = \{0\}$.
\end{proof}

\begin{theorem}[The linear extension exists and is unique]\label{thm:linear_ext}
    \lean{MPSTensor.linearExtension_exists_unique}
    \leanok
    \uses{def:injective, def:same_mpv}
    If $A$ is injective and $A$, $B$ generate the same matrix-product vectors,
    then there is a unique linear map $T : \MN{D} \to \MN{D}$ with $T(A^i) = B^i$
    for all~$i$.
\end{theorem}

\begin{proof}
    \leanok
    \uses{thm:trace_pairing_injective, lem:same_mpv_trace, lem:injectivity_transfer}
    Length-two trace agreement gives $\Phi_A(A^i) = \Phi_B(B^i)$ for every~$i$.
    The $\{A^i\}$ span $\MN{D}$, so $\operatorname{range} \Phi_A \subseteq
        \operatorname{range} \Phi_B$. With $\ker \Phi_A = \{0\}$
    (Theorem~\ref{thm:trace_pairing_injective}),
    Lemma~\ref{lem:injectivity_transfer} gives $\ker \Phi_B = \{0\}$. Let $g$ be
    a left inverse of~$\Phi_B$ and set $T = g \circ \Phi_A$; then
    \[
        T(A^i) = g(\Phi_A(A^i)) = g(\Phi_B(B^i)) = B^i.
    \]
    Uniqueness holds because the $\{A^i\}$ span $\MN{D}$.
\end{proof}

\begin{theorem}[The linear extension is multiplicative]\label{thm:linear_ext_mul}
    \lean{MPSTensor.linearExtension_mul}
    \leanok
    \uses{def:injective, def:same_mpv}
    Under the hypotheses of Theorem~\ref{thm:linear_ext}, the map~$T$ satisfies
    $T(MN) = T(M) T(N)$ for all $M, N \in \MN{D}$.
\end{theorem}

\begin{proof}
    \leanok
    \uses{thm:trace_pairing_injective, lem:same_mpv_trace, lem:injectivity_transfer}
    Length-two trace agreement gives $\Phi_B(T(A^i)) = \Phi_B(B^i) =
        \Phi_A(A^i)$; since the $\{A^i\}$ span $\MN{D}$ this extends to
    $\Phi_B \circ T = \Phi_A$, so $\operatorname{range} \Phi_A =
        \operatorname{range}(\Phi_B \circ T) \subseteq \operatorname{range} \Phi_B$. With
    $\ker \Phi_A = \{0\}$ (Theorem~\ref{thm:trace_pairing_injective}),
    Lemma~\ref{lem:injectivity_transfer} makes $\Phi_B$ injective. Length-three
    trace agreement gives, for all $i, j, k$,
    \[
        \Phi_B(T(A^i A^j))_k = \Phi_A(A^i A^j)_k = \tr(A^i A^j A^k)
        = \tr(B^i B^j B^k) = \Phi_B(B^i B^j)_k,
    \]
    so injectivity of $\Phi_B$ yields $T(A^i A^j) = B^i B^j$. Extend in two
    stages: for fixed~$j$, the maps $M \mapsto T(M A^j)$ and $M \mapsto T(M) B^j$
    agree on the spanning $\{A^i\}$, hence $T(M A^j) = T(M) B^j$ for all~$M$;
    then for fixed~$M$, the maps $N \mapsto T(M N)$ and $N \mapsto T(M) T(N)$
    agree on the spanning $\{A^j\}$, hence $T(M N) = T(M) T(N)$ for all~$N$.
\end{proof}

\section{Inner automorphism and the single-block theorem}

The rigidity of the full matrix algebra now finishes the proof. A nonzero
multiplicative map on~$\MN{D}$ is forced to be bijective (simplicity),
unit-preserving (surjectivity), and conjugation by a single invertible matrix
(Skolem--Noether). The next four results establish these steps and the
nonvanishing of~$T$, which the final theorem combines into the gauge~$X$.

\begin{theorem}[A nonzero multiplicative endomorphism of $\MN{D}$ is bijective]%
    \label{thm:simplicity}
    \lean{MPSTensor.linear_mul_endomorphism_bijective}
    \leanok
    A nonzero multiplicative $\C$-linear endomorphism of $\MN{D}$ is bijective.
\end{theorem}

\begin{proof}
    \leanok
    The kernel of a multiplicative linear map is a two-sided ideal. The algebra
    $\MN{D}$ is simple, so the kernel is $\{0\}$ or all of $\MN{D}$; the latter
    forces the map to be zero. Hence the map is injective, and surjective by
    finite-dimensionality.
\end{proof}

\begin{lemma}[A surjective multiplicative map preserves the unit]\label{lem:linear_map_to_alg}
    \lean{MPSTensor.linearMapToAlgHom}
    \leanok
    If $T : \MN{D} \to \MN{D}$ is a surjective multiplicative $\C$-linear map,
    then $T(\Id) = \Id$, so $T$ is a $\C$-algebra homomorphism.
\end{lemma}

\begin{proof}
    \leanok
    Choose~$X$ with $T(X) = \Id$. Then $T(\Id) = T(X) T(\Id) = T(X \Id) = T(X) =
        \Id$.
\end{proof}

\begin{lemma}[An injective tensor forces a nonzero partner]\label{lem:trace_ne_zero}
    \lean{MPSTensor.trace_ne_zero_of_injective}
    \leanok
    \uses{def:injective, def:same_mpv}
    Let $D \ge 1$. If $A$ is injective and $A$, $B$ generate the same
    matrix-product vectors, then the $B^i$ are not all zero.
\end{lemma}

\begin{proof}
    \leanok
    \uses{lem:same_mpv_trace}
    If $B^i = 0$ for every~$i$, then Lemma~\ref{lem:same_mpv_trace} on
    length-one words gives $\tr(A^i) = 0$ for all~$i$. The $\{A^i\}$ span
    $\MN{D}$, so the trace functional vanishes identically, contradicting
    $\tr(\Id) = D \neq 0$.
\end{proof}

\begin{theorem}[Skolem--Noether]\label{thm:skolem_noether}
    \lean{MPSTensor.skolemNoether_matrix}
    \leanok
    Every $\C$-algebra automorphism $f$ of $\MN{D}$ is inner: there is $X \in
        \GL_{D}(\C)$ with $f(M) = X M X^{-1}$ for all~$M$.
\end{theorem}

\begin{proof}
    \leanok
    Under the identification $\MN{D} \cong \operatorname{End}_{\C}(\C^D)$, every
    algebra automorphism of $\operatorname{End}_{\C}(\C^D)$ is conjugation by an
    invertible linear map. Translating back gives the invertible~$X$.
\end{proof}

\begin{theorem}[Single-block Fundamental Theorem]\label{thm:ft_single}
    \lean{MPSTensor.fundamentalTheorem_singleBlock}
    \leanok
    \uses{def:injective, def:same_mpv, def:gauge_equiv}
    Let $A$ be an injective MPS tensor and $B$ a tensor of the same bond
    dimension. If $A$ and $B$ generate the same matrix-product vectors, then
    there is $X \in \GL_{D}(\C)$ with $B^i = X A^i X^{-1}$ for all~$i$.
\end{theorem}

\begin{proof}
    \leanok
    \uses{thm:linear_ext, thm:linear_ext_mul, thm:simplicity,
        lem:trace_ne_zero, lem:linear_map_to_alg, thm:skolem_noether}
    Theorem~\ref{thm:linear_ext} gives the unique linear~$T$ with $T(A^i) = B^i$,
    and Theorem~\ref{thm:linear_ext_mul} makes it multiplicative. It is nonzero:
    otherwise $B^i = 0$ for all~$i$, against Lemma~\ref{lem:trace_ne_zero}. By
    Theorem~\ref{thm:simplicity}, $T$ is bijective, and by
    Lemma~\ref{lem:linear_map_to_alg} it fixes~$\Id$, so $T$ is a $\C$-algebra
    automorphism. Theorem~\ref{thm:skolem_noether} gives $X \in \GL_{D}(\C)$ with
    $T(M) = X M X^{-1}$; evaluating at $M = A^i$ gives $B^i = X A^i X^{-1}$.
\end{proof}